% ============================================================
%  Chương IV – Chứng chỉ số và các chuẩn
%  ~900 từ · 2.5 trang
% ============================================================
\chapter{Chứng chỉ số và các chuẩn}
\label{chap:x509}

% Ý chính: chứng chỉ số là "sản phẩm" cốt lõi của PKI.
% Chương đi sâu vào format, cấu trúc nội tại, phân loại,
% và hệ sinh thái chuẩn hoá xung quanh.

% ----------------------------------------------------------
\section{Chứng chỉ số X.509}
\label{sec:4.1}

\subsubsection*{Lịch sử và phiên bản}
% Ý chính: v1 (1988, ITU-T X.500) → v2 (1993, thêm UID) → v3 (1996, extensions) \cite{RFC5280}
% Định nghĩa trực quan: cert = CA ký lên (identity + public key)

\subsubsection*{Cấu trúc X.509 v3}
% Ý chính: Version, SerialNumber, SignatureAlgorithm, Issuer,
% Validity (notBefore/notAfter), Subject, SubjectPublicKeyInfo, Extensions, Signature
% Bảng~\ref{tab:x509-fields}: liệt kê và giải thích từng trường

\subsubsection*{Distinguished Name (DN)}
% Ý chính: C, ST, L, O, OU, CN — cách nhận dạng entity trong X.500 namespace
% Lưu ý: CN khác SAN; CN không dùng để validate domain kể từ RFC 2818

\subsubsection*{Extensions quan trọng}
% Ý chính: basicConstraints (isCA + pathLenConstraint), keyUsage, extendedKeyUsage,
% Subject Alternative Name (SAN) — trường thực sự dùng để validate domain,
% CRL Distribution Point, Authority Information Access (OCSP URL)

% ----------------------------------------------------------
\section{Certificate Signing Request (CSR)}
\label{sec:4.2}
% Ý chính: entity tự tạo keypair, tự ký CSR bằng private key (proof of possession),
% gửi CSR kèm thông tin danh tính cho CA xác minh rồi ký
% Format: PKCS\#10 (RFC 2986); các trường DN + public key + signature

% ----------------------------------------------------------
\section{Phân loại chứng chỉ}
\label{sec:4.3}

\subsubsection*{Theo mức độ xác minh danh tính}
% Bảng~\ref{tab:cert-types}: DV / OV / EV
% DV: chỉ verify domain control (DNS/HTTP challenge); phát hành tự động trong phút
% OV: verify tổ chức qua hồ sơ pháp lý; EV: mức cao nhất, dùng code signing

\subsubsection*{Theo phạm vi tên miền}
% Single domain: chỉ khớp một FQDN
% Wildcard (*.example.com): khớp một cấp subdomain, không khớp sub-subdomain
% Multi-SAN: nhiều domain hoàn toàn khác nhau trong một cert

% ----------------------------------------------------------
\section{Các chuẩn PKCS}
\label{sec:4.4}
% Bảng tóm tắt họ chuẩn PKCS của RSA Laboratories:
%   PKCS\#10 — định dạng CSR
%   PKCS\#12 (PFX) — đóng gói cert + private key (dùng để export/import)
%   PKCS\#11 (Cryptoki) — interface lập trình với HSM và smart card
%   PKCS\#7 / CMS — định dạng dữ liệu ký số và mã hóa (dùng trong S/MIME)

% ----------------------------------------------------------
\section{RFC và chuẩn quan trọng}
\label{sec:4.5}
% Ý chính: ecosystem chuẩn hoá xung quanh PKI
%   RFC 5280 \cite{RFC5280} — X.509 cert \& CRL profile (nền tảng)
%   RFC 6960 \cite{RFC6960} — OCSP
%   RFC 8446 \cite{RFC8446} — TLS 1.3
%   RFC 8555 \cite{RFC8555} — ACME Protocol
%   CA/Browser Forum Baseline Requirements \cite{CABForum} — quy tắc vận hành CA
